\section{Special-function expansions and absolute error envelopes}
\label{sec:special-function-expansions}

The preceding calculus does not depend on an elementary description of
the function being expanded. Its analytic input is a finite approximation
on a specified branch, together with a bound for the omitted part.
Special functions supply such input through differential equations,
integral representations, and convergent defining sums. These sources
must be distinguished from exact identities: a finite asymptotic
expansion need not determine contributions smaller than every power in
its stated scale.

All parameters in this section are fixed while the argument approaches
its endpoint. Unless an explicit bound is given, constants and
thresholds in $\OO$ may depend on those parameters and on the fixed
truncation order. No uniformity in a growing order, a coalescing
singularity, or a parameter approaching the boundary of its domain is
inferred. We use the real positive axis for fractional powers and
logarithms. Reflection or continuation to another ray requires an exact
identity or a separate asymptotic theorem on that ray.

\subsection{Finite composition, oscillation, and real projection}

Let $w\to0^+$, and suppose finitely many functions have representations
\begin{equation}
 F_j(w)=A_j(w)+\OO(R_j(w)),\qquad R_j(w)\ge0.
 \label{eq:special-function-input-envelope}
\end{equation}
The approximations may contain exact exponential factors and bounded
oscillations. For example, a finite expression of the form
\begin{equation}
 A(w)=A_0(w)+\sum_{j=1}^m P_j(w)T_j(w)
 \label{eq:special-function-finite-carriers}
\end{equation}
can keep the factors $P_j$ exact while each $T_j$ is a finite
power--logarithmic sum with trigonometric coefficients. An error
$\OO(R_j)$ inside the $j$th bracket gives the absolute error
$\OO(|P_j|R_j)$ outside it. Several independent errors give their sum
as a valid envelope, even when no common leading term is available.

Proposition~\ref{prop:separate-error-envelopes} propagates these errors
through finite arithmetic expressions. At a regular finite limiting
argument, Taylor's theorem supplies composition with an analytic
function. More generally, a derivative bound on a neighborhood containing
the true and approximate arguments supplies the corresponding
mean-value bound. Division requires separation from zero; a real
noninteger power requires the appropriate positive branch; exponentiation
of an uncertain phase requires its absolute error to tend to zero.
Finitely many successive applications are justified by induction, with
every earlier error included at the next step.

An oscillatory coefficient changes the interpretation of the result.
For instance,
\[
 F(x)=x^{-1/2}\cos x+\OO(x^{-3/2})
\]
is a useful absolute approximation on the positive axis. It does not
assert $F(x)/(x^{-1/2}\cos x)\to1$, since the denominator vanishes
arbitrarily far along that axis. Nor does it justify taking a reciprocal
on a whole positive half-line. The boundedness of sine and cosine
supports absolute error transport through their zeros.

\begin{proposition}[Projection of a real-valued function]
\label{prop:special-function-real-projection}
Suppose $F$ is real-valued on the specified approach, $A$ may be complex,
and $F-A=\OO(R)$ for a nonnegative real envelope $R$. Then
\[
 F-\operatorname{Re}A=\OO(R),
 \qquad \operatorname{Im}A=\OO(R).
\]
In particular, projecting the finite approximation to its real part
does not increase an available pointwise absolute error bound.
\end{proposition}
\begin{proof}
Since $\operatorname{Re}F=F$ and $\operatorname{Im}F=0$, both assertions
follow from $|\operatorname{Re}z|\le|z|$ and
$|\operatorname{Im}z|\le|z|$. The same inequalities hold before an
explicit bound is replaced by asymptotic notation.
\end{proof}

The real-valuedness assumption refers to the exact function on its
chosen branch. It cannot be deduced merely by discarding an imaginary
part of an approximation. Likewise, algebraic cancellation of finite
terms need not cancel their unknown errors. If
$F=A+B+\OO(R)$, $B=\OO(S)$ and $R=\oo(S)$, then one may conclude
$F=A+\OO(S)$; the smaller independent error has been absorbed into a
proved bound for the omitted expression.

\subsection{An exponent allowance for logarithmic tails}

A remainder exponent without a logarithmic degree is incomplete
information in the scale of Section~\ref{sec:scale}. If it is known
independently that the degree is finite, a strict loss in exponent gives
a bound that does not need that degree explicitly.

\begin{lemma}[Absorbing an unspecified finite logarithmic degree]
\label{lem:special-function-logarithmic-margin}
Suppose, for some fixed real $p$ and finite $D\ge0$,
\[
 R(w)=\OO\bigl(w^p\M(w)^D\bigr).
\]
For every fixed $\varepsilon>0$,
\[
 R(w)=\oo\bigl(w^{p-\varepsilon}\bigr).
\]
In particular, $R(w)=\OO(w^{p-\varepsilon})$.
\end{lemma}
\begin{proof}
Divide the assumed bound by $w^{p-\varepsilon}$ and apply
Lemma~\ref{lem:dominance} to $w^\varepsilon\M(w)^D$.
\end{proof}

For a lattice of exponent spacing $1/q$, choosing
$\varepsilon=1/(2q)$ leaves an intermediate exponent between adjacent
lattice positions. This choice is only a conservative bound; it does not
specify the actual first omitted coefficient. A subsequently established
omitted block may give a sharper boundary by the absorption argument
above. The finite-degree hypothesis is essential: a purely formal
remainder symbol or the logarithmic degrees of retained coefficients do
not prove it. These assertions concern values, and provide no derivative
bounds without the additional hypotheses in
Section~\ref{sec:series-calculus}; see also \cite[\S2.1(iii)]{dlmf}.

More precisely, suppose a complete omitted block has been determined in
an approximation
\[
 F(w)=T(w)+w^\beta C(\log w)+E(w),\qquad
 E(w)=\OO(w^\rho\M(w)^D),\qquad \rho>\beta,
\]
where $C$ is a nonzero polynomial and $D$ is fixed and finite.
Choosing $0<\varepsilon<\rho-\beta$ gives $E=\oo(w^\beta)$, and hence
$F-T=w^\beta C(\log w)+\oo(w^\beta)
=\OO(w^\beta\M^{\deg C})$. The strict gap is essential to this
argument: at $\rho=\beta$, the unknown degree $D$ cannot be replaced by
$\deg C$. Vanishing computed coefficients do not establish a zero
tail. The exponent allowance and any sharpening of it concern the
proved analytic bound, independently of the formal cutoff used to
select retained blocks. The constants and thresholds may depend on
fixed parameters; uniform estimates require corresponding uniform tail
hypotheses. No assertion here permits a logarithmic degree that varies
with $w$.

\subsection{Exact identities preserve smaller exponential terms}

On $x>0$, define the finite polynomials
\[
 P_n(t)=\sum_{k=0}^n
       \frac{(n+k)!}{k!(n-k)!2^k}\,t^k,
 \qquad n=0,1,2,\ldots.
\]
The half-integer modified Bessel functions satisfy
\begin{equation}
\begin{aligned}
 K_{n+1/2}(x)
   &=\sqrt{\frac{\pi}{2x}}e^{-x}P_n(1/x),\\
 I_{n+1/2}(x)
   &=\frac{e^xP_n(-1/x)+(-1)^{n+1}e^{-x}P_n(1/x)}
           {\sqrt{2\pi x}}.
\end{aligned}
\label{eq:special-function-half-integer-bessel}
\end{equation}
They follow from the half-order identities in
\cite[\S10.39]{dlmf} and the order recurrences. Here both equations
are exact. In particular,
\[
 e^{-x}I_{1/2}(x)
   =\frac{1-e^{-2x}}{\sqrt{2\pi x}}.
\]
Its inverse-power expansion contains only the first displayed algebraic
term, although the difference from that term is nonzero. The vanishing
of every subsequent algebraic coefficient therefore cannot establish
an exact finite identity.

The same distinction occurs for confluent hypergeometric functions:
\begin{equation}
 {}_1F_1(1;2;x)=\int_0^1e^{xt}\,\dd t
               =\frac{e^x-1}{x}\qquad(x\ne0).
 \label{eq:special-function-hypergeometric-endpoint}
\end{equation}
The value at zero is obtained by continuity. The integral is the
specialization of the Euler representation in
\cite[(13.4.1)]{dlmf}. Both endpoints contribute to the exact answer.
Keeping the term $-1/x$ matters after cancellation or rescaling, even
though it is exponentially small relative to $e^x/x$.

Terminating defining sums give a different route to exactness. If
$m\ge0$ is an integer and $b\notin\{0,-1,-2,\ldots\}$, then
\[
 {}_1F_1(-m;b;x)
   =\sum_{k=0}^{m}\frac{(-m)_k}{(b)_k\,k!}x^k.
\]
This conclusion comes from the zero numerator Pochhammer symbols and
the nonvanishing denominators in the defining series, rather than from
observing several zero coefficients in an asymptotic expansion.

\subsection{Exponential and oscillatory differential-equation families}

For fixed $\nu\in\R$, set
\[
 a_0(\nu)=1,\qquad
 a_k(\nu)=\frac{\prod_{j=1}^k(4\nu^2-(2j-1)^2)}{8^k k!},
 \qquad \omega=x-\frac{\nu\pi}{2}-\frac{\pi}{4}.
\]
The positive-axis modified Bessel expansions are, for fixed $N\ge1$,
\begin{equation}
\begin{aligned}
 I_\nu(x)
 &=\frac{e^x}{\sqrt{2\pi x}}
   \left(\sum_{k=0}^{N-1}\frac{(-1)^ka_k(\nu)}{x^k}
         +\OO(x^{-N})\right),\\
 K_\nu(x)
 &=\sqrt{\frac{\pi}{2x}}e^{-x}
   \left(\sum_{k=0}^{N-1}\frac{a_k(\nu)}{x^k}
         +\OO(x^{-N})\right).
\end{aligned}
\label{eq:special-function-modified-bessel}
\end{equation}
The corresponding oscillatory examples begin
\begin{equation}
\begin{aligned}
 J_\nu(x)&=\sqrt{\frac{2}{\pi x}}
   \left(\cos\omega-\frac{a_1(\nu)}x\sin\omega
       -\frac{a_2(\nu)}{x^2}\cos\omega+\OO(x^{-3})\right),\\
 Y_\nu(x)&=\sqrt{\frac{2}{\pi x}}
   \left(\sin\omega+\frac{a_1(\nu)}x\cos\omega
       -\frac{a_2(\nu)}{x^2}\sin\omega+\OO(x^{-3})\right).
\end{aligned}
\label{eq:special-function-oscillatory-bessel}
\end{equation}
These formulas and their finite-order bounds are given in
\cite[\S\S10.17, 10.40]{dlmf}. Each outer factor multiplies the
absolute error as well as the finite bracket. The oscillatory formulas
have this meaning throughout the positive axis, including zeros of
their leading trigonometric terms.

The Airy functions illustrate a fractional power in the exponential
phase. With $\xi=2x^{3/2}/3$ and $x\to+\infty$,
\begin{equation}
\begin{aligned}
 \operatorname{Ai}(x)
 &=\frac{e^{-\xi}}{2\sqrt{\pi}\,x^{1/4}}
   \left(1-\frac5{48x^{3/2}}+\frac{385}{4608x^3}
              +\OO(x^{-9/2})\right),\\
 \operatorname{Bi}(x)
 &=\frac{e^\xi}{\sqrt{\pi}\,x^{1/4}}
   \left(1+\frac5{48x^{3/2}}+\frac{385}{4608x^3}
              +\OO(x^{-9/2})\right).
\end{aligned}
\label{eq:special-function-airy-positive}
\end{equation}
On the reflected real ray they instead satisfy
\begin{equation}
\begin{aligned}
 \operatorname{Ai}(-x)&=\frac1{\sqrt{\pi}\,x^{1/4}}
        \left(\sin(\xi+\pi/4)+\OO(x^{-3/2})\right),\\
 \operatorname{Bi}(-x)&=\frac1{\sqrt{\pi}\,x^{1/4}}
        \left(\cos(\xi+\pi/4)+\OO(x^{-3/2})\right).
\end{aligned}
\label{eq:special-function-airy-negative}
\end{equation}
See \cite[\S9.7]{dlmf}. Thus a change of ray can replace real
exponential growth by oscillation. The derivative expansions in that
reference are separate analytic assertions; the displayed value
remainders alone would not justify differentiating them.

Exact substitutions such as $x=c/w^\alpha+d$, with $c,\alpha>0$ and
fixed real $d$, preserve a positive-axis asymptotic theorem once the
argument is eventually positive. The remainder is first substituted
in its full absolute form. Further expansion of the exact phase or
prefactor then uses the finite-composition rules, with enough accuracy
in a phase to control its exponential.

\subsection{A finite-moment theorem for integral and sum representations}

A single Taylor estimate treats several integral tails and discrete
defining sums. Write $|\mu|$ for the total variation of a real signed
measure $\mu$ on $[0,\infty)$, and
$M_k(\mu)=\int t^k\,\dd\mu(t)$ for its moments.

\begin{theorem}[Expansion against a fixed measure]
\label{thm:special-function-moment-expansion}
Fix $s\in\R$ and an integer $N\ge0$, and set
$D=\max(0,\lceil-s-N\rceil)$. Suppose
\[
 \int_0^\infty(1+t)^{N+D}\,\dd|\mu|(t)<\infty.
\]
For $a\ge1$, define
\[
 F_\mu(s,a)=\int_0^\infty(a+t)^{-s}\,\dd\mu(t).
\]
Then
\begin{equation}
 F_\mu(s,a)=
 \sum_{k=0}^{N-1}\frac{(-1)^k(s)_kM_k(\mu)}{k!}\,a^{-s-k}
       +R_N(a),
 \label{eq:special-function-moment-expansion}
\end{equation}
with the empty sum understood as zero, and
\begin{equation}
\begin{aligned}
 |R_N(a)|&\le C_N(\mu,s)\,a^{-s-N},\\
 C_N(\mu,s)&=
  \frac{|(s)_N|}{N!}
       \int_0^\infty t^N(1+t)^D\,\dd|\mu|(t).
\end{aligned}
\label{eq:special-function-moment-bound}
\end{equation}
\end{theorem}
\begin{proof}
For $h(v)=(1+v)^{-s}$ on $v\ge0$,
$h^{(N)}(v)=(-1)^N(s)_N(1+v)^{-s-N}$. If $N\ge1$,
Taylor's remainder on $[0,v]$ is therefore bounded by
\[
 \left|h(v)-\sum_{k=0}^{N-1}
                \frac{(-1)^k(s)_k}{k!}v^k\right|
       \le \frac{|(s)_N|}{N!}\,v^N(1+v)^D.
\]
For $N=0$ the same inequality is simply
$(1+v)^{-s}\le(1+v)^D$.
Substitute $v=t/a$, multiply by $a^{-s}$, and use
$(1+t/a)^D\le(1+t)^D$ for $a\ge1$. Integrating the absolute
bound against $|\mu|$ proves both formulas. The moment hypothesis
ensures that the original integral, every retained term, and the
remainder bound are finite.
\end{proof}

If all moments of $|\mu|$ are finite, the theorem gives every fixed
order of a Poincar\'e expansion. It does not assert convergence as
$N\to\infty$. If $s=-m$ for an integer $m\ge0$, the binomial expression
is a polynomial and the expansion is exact after degree $m$.

\paragraph{Exponential integrals and incomplete Gamma functions.}
For fixed real $p$ and $a>0$, the positive-axis integral
\[
 E_p(a)=\int_1^\infty e^{-at}t^{-p}\,\dd t
       =\frac{e^{-a}}a\int_0^\infty
                  e^{-v}(1+v/a)^{-p}\,\dd v
\]
has moments $\int_0^\infty e^{-v}v^k\,\dd v=k!$.
Theorem~\ref{thm:special-function-moment-expansion} gives
\begin{equation}
 E_p(a)=\frac{e^{-a}}a
       \left(\sum_{k=0}^{N-1}\frac{(-1)^k(p)_k}{a^k}
                   +\OO(a^{-N})\right).
 \label{eq:special-function-exponential-integral}
\end{equation}
When $p+N\ge0$, the absolute remainder inside the bracket is at most
$|(p)_N|a^{-N}$ for $a\ge1$.
The identities
\[
 \Gamma(\alpha,a)=a^\alpha E_{1-\alpha}(a),\qquad
 \operatorname{erfc}x=\frac{\Gamma(1/2,x^2)}{\sqrt{\pi}}
       \quad(x>0)
\]
give incomplete-Gamma and complementary-error-function tails by the
same argument. Their standard forms and sharper domains for bounds
are recorded in \cite[\S\S8.11(i), 7.12(i)]{dlmf}.
Section~\ref{sec:special-function-adapters} establishes the additional
derivative bounds needed for complementary-error-function inversion.

The growing exponential integral has a different positive-axis form,
\begin{equation}
 \operatorname{Ei}(x)=\frac{e^x}{x}
       \left(\sum_{k=0}^{N-1}\frac{k!}{x^k}+\OO(x^{-N})\right).
 \label{eq:special-function-growing-exponential-integral}
\end{equation}
The real reflection identity
$\operatorname{Ei}(-x)=-E_1(x)$ for $x>0$ selects its negative-axis
tail. The sine and cosine integrals instead begin
\begin{equation}
\begin{aligned}
 \operatorname{Si}(x)&=\frac{\pi}{2}-\frac{\cos x}{x}
                              -\frac{\sin x}{x^2}+\OO(x^{-3}),\\
 \operatorname{Ci}(x)&=\frac{\sin x}{x}
                              -\frac{\cos x}{x^2}+\OO(x^{-3}).
\end{aligned}
\label{eq:special-function-oscillatory-integrals}
\end{equation}
These positive-axis formulas are in \cite[\S6.12]{dlmf}.
They again give absolute approximations through oscillatory zeros.

\paragraph{A confluent hypergeometric example.}
For fixed $b\in\R$, $\alpha>0$ and $x>0$, the integral representation
\cite[(13.4.4)]{dlmf}, after rescaling its integration variable, gives
\[
 U(\alpha,b,x)=\frac{x^{-\alpha}}{\Gamma(\alpha)}
     \int_0^\infty e^{-t}t^{\alpha-1}
                      (1+t/x)^{-(1+\alpha-b)}\,\dd t.
\]
The positive measure here has moments
$M_k=(\alpha)_k$. Taking $s=1+\alpha-b$ in the moment theorem yields
\begin{equation}
 U(\alpha,b,x)=x^{-\alpha}
   \left(\sum_{k=0}^{N-1}
        \frac{(-1)^k(\alpha)_k(1+\alpha-b)_k}{k!\,x^k}
          +\OO(x^{-N})\right).
 \label{eq:special-function-tricomi-moments}
\end{equation}
The moment bound supplies an explicit constant at each fixed order.
Extending this proof to other values of $\alpha$ requires a
continuation or recurrence argument with its own domain conditions.

\subsection{Lerch's transcendent with a large third argument}

For fixed real $z,s$ with $|z|<1$ and $a>0$, the defining sum is
\begin{equation}
 \Phi(z,s,a)=\sum_{n=0}^\infty z^n(a+n)^{-s}.
 \label{eq:special-function-lerch-definition}
\end{equation}
It converges absolutely for every real $s$; all powers use positive
bases. See \cite[(25.14.1)]{dlmf}. The signed discrete measure
$\mu_z=\sum_{n\ge0}z^n\delta_n$ has finite absolute moments of every
order. Define
\begin{equation}
 M_0(z)=\frac1{1-z},\qquad
 M_k(z)=\sum_{n=1}^\infty n^kz^n
       =\operatorname{Li}_{-k}(z)\quad(k\ge1).
 \label{eq:special-function-lerch-moments}
\end{equation}
The last equality is the polylogarithm defining sum
\cite[(25.12.10)]{dlmf}; the separate zeroth moment includes the
mass at $n=0$.

\begin{corollary}[Lerch expansion and a computable tail constant]
\label{cor:special-function-lerch}
For every fixed integer $N\ge0$,
\begin{equation}
 \Phi(z,s,a)=\sum_{k=0}^{N-1}
       \frac{(-1)^k(s)_k M_k(z)}{k!}\,a^{-s-k}
              +R_N(z,s,a),
 \label{eq:special-function-lerch-expansion}
\end{equation}
and, for $a\ge1$,
\begin{equation}
\begin{aligned}
 |R_N(z,s,a)|&\le C_N(z,s)\,a^{-s-N},\\
 C_N(z,s)&=\frac{|(s)_N|}{N!}
       \sum_{j=0}^{D}\binom Dj M_{N+j}(|z|),
 \qquad D=\max(0,\lceil-s-N\rceil).
\end{aligned}
\label{eq:special-function-lerch-bound}
\end{equation}
\end{corollary}
\begin{proof}
Apply Theorem~\ref{thm:special-function-moment-expansion} to $\mu_z$.
Its total variation is $\mu_{|z|}$. Expanding the finite polynomial
$(1+n)^D$ gives the stated constant.
\end{proof}

Negative values of $z$ change the coefficients through cancellation;
the bound uses $|z|$ and remains valid. Negative noninteger values of
$s$ give growing initial powers and are covered by the factor
$(1+t)^D$ in the proof. For example,
\begin{equation}
\begin{aligned}
 \Phi(1/2,2,a)
   &=\frac2{a^2}-\frac4{a^3}+\frac{18}{a^4}
                   +\OO(a^{-5}),\\
 \Phi(1/2,-1/2,a)
   &=2\sqrt a+\frac1{\sqrt a}+\OO(a^{-3/2}).
\end{aligned}
\label{eq:special-function-lerch-examples}
\end{equation}
The respective absolute remainder constants from
\eqref{eq:special-function-lerch-bound} are $104$ and $3/4$ for
$a\ge1$. If $z=0$, the exact answer is $a^{-s}$. If $s=-m$ with
$m\ge0$ an integer, the expansion terminates:
\[
 \Phi(z,-m,a)=\sum_{k=0}^m\binom mk a^{m-k}M_k(z).
\]
For example $\Phi(1/2,-2,a)=2a^2+4a+6$.

The positive coordinate is $w=1/a$, with absolute block exponents
$s+k$. Zero moments or zero Pochhammer factors remove entire blocks;
they do not authorize changing the scale of the remaining uncertainty.
A strict cutoff $H$ retains only nonzero blocks with $s+k<H$.
The constants above need not stay bounded as $|z|\to1$, and the
finite-measure proof does not apply at $z=\pm1$. At $z=1$, for example,
the function becomes the Hurwitz zeta function when $s>1$, and integral
comparison gives leading behavior $a^{1-s}/(s-1)$, a different power
from the fixed-$|z|<1$ leading term $a^{-s}/(1-z)$.

\subsection{The zeta function in an exponential coordinate}

Large positive values of the argument of the Riemann zeta function
lead to a different scale. Its convergent defining sum is
\cite[(25.2.1)]{dlmf}
\[
 \zeta(S)=\sum_{n=1}^\infty n^{-S},\qquad S>1.
\]
With $w=e^{-S}\to0^+$ this becomes
\begin{equation}
 \zeta(-\log w)=\sum_{n=1}^\infty w^{\log n},
       \qquad 0<w<e^{-1}.
 \label{eq:special-function-zeta-coordinate}
\end{equation}
The exponents $\{\log n:n\ge1\}$ are left-finite, since
$\log n<H$ implies $n<e^H$. They form an additive semigroup under
$\log m+\log n=\log(mn)$, and are not finitely generated: a finite
set of integers involves only finitely many prime factors.
Every finite cutoff still has finitely many
terms, and multiplication has finite coefficient sums indexed by
factorizations of an integer. Thus the finite arithmetic extends
beyond a globally finitely generated exponent semigroup.

\begin{proposition}[A two-sided zeta tail bound]
\label{prop:special-function-zeta-tail}
For every integer $m\ge1$ and every real $S>1$,
\begin{equation}
\begin{aligned}
 \zeta(S)&=\sum_{n=1}^{m-1}n^{-S}+R_m(S),\\
 m^{-S}&\le R_m(S)
       \le m^{-S}\left(1+\frac{m}{S-1}\right).
\end{aligned}
\label{eq:special-function-zeta-tail}
\end{equation}
Consequently, for fixed $m$, $R_m(S)\sim m^{-S}$ as $S\to+\infty$.
\end{proposition}
\begin{proof}
Every summand of $R_m$ is positive, giving the lower bound. Since
$t^{-S}$ decreases on $t>0$,
\[
 \sum_{n=m+1}^\infty n^{-S}
       \le\int_m^\infty t^{-S}\,\dd t
       =\frac{m^{1-S}}{S-1}.
\]
Add $m^{-S}$ and divide both bounds by it. The quotient of the
remainder and its first term is squeezed to $1$.
\end{proof}

For example, retaining five terms gives
\begin{equation}
 \zeta(S)=1+2^{-S}+3^{-S}+4^{-S}+5^{-S}+\OO(6^{-S}).
 \label{eq:special-function-zeta-five-terms}
\end{equation}
Here the first term is the constant corresponding to exponent
$\log1=0$. A strict cutoff in the coordinate $w=e^{-S}$ retains
$\log n<H$ and has first omitted exponent $\log m$ as in
\eqref{eq:special-function-zeta-tail}. Substituting
$S=ax+b\to+\infty$ preserves both bounds on $ax+b>1$.
In contrast, all corrections to $1$ are flat in $1/S$; an
inverse-power expansion in that coordinate cannot distinguish them.

The zeta representation is a convergent generalized power series on
its stated interval. The Lerch expansion follows from a convergent
defining sum but is asserted only at each fixed asymptotic order.
This distinction persists under finite composition, together with the
absolute envelopes, branch conditions, and fixed-parameter hypotheses
used at each step.
